\section{Results, source analysis, and scope}\label{sec:context}

The main result is an exact relation between two objects: the zeros of a
polynomial constructed from a graph, and the electrical resistances of that
graph.  The relation can be computed before any of the input zero positions
are specified.  It also controls what happens when a single graph edge is
strengthened.  Theorem~\ref{thm:electrical} gives the full statement and
proof, and Theorem~\ref{thm:edge-law} gives the edge update.

The arithmetic application starts from the established Dirichlet
zero-location theorem in the supplied Lean repository.  A completed
$L$-function is rotated so that its critical line becomes the real axis.
The graph construction then gives further entire functions with all zeros
on that axis.  The finite electrical formulas apply to every finite packet
of zeros of the input functions.  The entire-function construction has its
own direct proof, so it does not require summing infinitely many zero
displacements.

\subsection{What the supplied sources actually establish}

The source inspection used the following fixed revisions:
\begin{center}
\small
\begin{tabular}{p{.26\textwidth}p{.66\textwidth}}
\toprule
Source & Revision or retrieval\\
\midrule
\texttt{lsernau/grh-lean} &
\path{72fead102f663d5c7a407964a13dc0cf9d977207}\\
\texttt{ProveIt} &
\path{4d7990202f9f106b2ffe8c209f073ad01d5e1987}\\
Wikipedia, Riemann hypothesis & Retrieved 26 November 2026; used for context\\
\bottomrule
\end{tabular}
\end{center}

The statement module of the compact Lean edition imports Mathlib and
formulates, among its endpoints, nonvanishing for every Dirichlet character
$\chi$ modulo every positive integer $q$:
\[
 \Re s>\tfrac12,\qquad (\chi\ne1\ \text{or}\ s\ne1)
 \quad\Longrightarrow\quad L(s,\chi)\ne0.
\]
It separately states the critical-strip zero-location theorem.  The
primitive-character version is a specialization, while the principal
character gives the zeta endpoint after the usual treatment of the pole.
The source does not state GRH for all Artin, Hecke, or automorphic
$L$-functions.\cite{grhsource}

The compact guide describes a comparison of covolumes obtained through
chains of theta lifts: sufficiently many short forms are compared with a
lower bound for integral forms.  This produces a one-sided estimate for a
smoothed prime remainder.  The final analytic step uses meromorphic
Laplace transforms, a quadratic twist that prevents residue cancellation,
and a Landau-type one-sign argument.  This is the mechanism visible in
the inspected guide and endgame modules.  The present graph argument is a
subsequent use of the zero-location endpoint and determinant methods; it
is not an alternative explanation of that theta-lift proof.\cite{grhguide}

The upstream README reports machine verification with the standard Lean
axioms \texttt{propext}, \texttt{Classical.choice}, and \texttt{Quot.sound},
and says the development has not been read through by a mathematician.
The Lean FRO and Mathlib Initiative's separate public statement endorses
Alp\"oge's original RH proof following statement comparison and independent
proof checks.  The retrieved Wikipedia article reports that development.
Those are externally reported checks; this packet makes no claim to have
reproduced the complete upstream build.\cite{grhsource,leanstatement,wikipedia}

\subsection{Comparison with ProveIt}

The organized GRH directory contains 61 research packets and explicitly
labels their placement as unreviewed.  At the pinned revision its local
\texttt{docs/incoming} directory also contains 39 later ZIP packets.  We
read all 61 organized delivery READMEs, extracted the text members of the
39 later archives, inspected their delivery summaries, and searched the
resulting corpus for the proposed construction.  This is a comparison of
the reported results and selected text, not an independent verification of
100 papers.  The exact inventory is recorded in the accompanying
provenance files.\cite{proveit}

Many plausible first directions are already represented there: Li
coefficients, Hankel and Ritz recovery of zeros, finite noisy certificates,
sampling, filtered prime races, torsion specialization, graph-tower
valuations, and symplectic or orthogonal matrix blocks.  The present
construction concerns a different graph operation: a rooted-forest sum of
functions and derivatives, an explicit symmetric compression, and an
electrical formula for the resulting zero displacement.  No matching
statement was found in the inspected corpus.

\subsection{What is classical and what is proposed as new}

The following ingredients are established mathematics: the matrix-forest
identity; compression formulas representing the critical points of a
polynomial; preservation of real zeros in the theory of stable
polynomials; matrix-valued Herglotz theory; and rank-one matrix inverse
identities.  Relevant primary sources include Chebotarev--Shamis,
Pereira, Borcea--Br\"and\'en, Reiffenstein, and
Simonov--Woracek.\cite{forests,pereira,bb,reiffenstein,simonov}

The contribution proposed here combines the rooted-forest derivative sum
with the explicit symmetric compression
\[
 K_G=(Q\Lambda Q-QVL^+V^{\mathsf T}Q)|_{\boldsymbol1_N^\perp}.
\]
From this we derive the exact resistance formula for zero transport,
its exact single-edge displacement law, and the completed
Dirichlet $L$-function specialization.  The proofs below establish the
claims as mathematical statements.  A targeted search did not locate this
precise electrical transport formula, but that is not a proof of first
publication.  In particular, the basic real-zero preservation theorem is
not presented as a newly discovered general principle.

The polynomial theorems are unconditional.  The arithmetic specialization
uses the supplied Dirichlet GRH result solely to put the completed
$L$-function zeros on the real axis.  It yields explicit differential
combinations; it does not identify new Euler products or prove a further
automorphic GRH.
